% =============================================================================
% TEMPLATE DE ARTIGO - PROJETO DE PESQUISA EM COMPUTAÇÃO GRÁFICA
% Bacharelado em Ciência da Computação - IFNMG - Campus Montes Claros
% Disciplina: Computação Gráfica - 2026/2
% Professor: Wagner Ferreira de Barros
%
% Formato didático inspirado em características recorrentes de conferências de
% Computação Gráfica (p.ex., ACM SIGGRAPH/SIGGRAPH Asia e Eurographics):
% texto compacto, duas colunas, corpo em 9 pt, pequeno espaçamento vertical e
% figuras/tabelas integradas ao fluxo.
%
% IMPORTANTE: este é o modelo oficial DA DISCIPLINA, não o template oficial de
% um evento específico. Para futura submissão, migre o conteúdo para o template
% oficial do evento escolhido e siga integralmente as regras daquela chamada.
%
% Extensão exigida na disciplina: 8 a 12 páginas, incluindo referências.
% =============================================================================

\documentclass[9pt,a4paper,twocolumn]{extarticle}

% -----------------------------------------------------------------------------
% Codificação, idioma e tipografia
% -----------------------------------------------------------------------------
\usepackage[T1]{fontenc}
\usepackage[utf8]{inputenc}
\usepackage[brazil]{babel}
\usepackage{newtxtext,newtxmath}
\usepackage{microtype}

% -----------------------------------------------------------------------------
% Página e colunas
% Área útil próxima aos padrões compactos usados em publicações de CG.
% -----------------------------------------------------------------------------
\usepackage[a4paper,textwidth=177mm,textheight=245mm,centering]{geometry}
\setlength{\columnsep}{8mm}
\setlength{\columnseprule}{0pt}
\setlength{\parindent}{1em}
\setlength{\parskip}{0pt}
\linespread{0.97}
\pagestyle{empty}

% -----------------------------------------------------------------------------
% Seções compactas
% -----------------------------------------------------------------------------
\usepackage{titlesec}
\titleformat{\section}
  {\normalfont\bfseries\fontsize{10.5}{11.5}\selectfont}
  {\thesection}{0.5em}{}
\titleformat{\subsection}
  {\normalfont\bfseries\fontsize{9.5}{10.5}\selectfont}
  {\thesubsection}{0.45em}{}
\titlespacing*{\section}{0pt}{7pt plus 2pt minus 1pt}{3pt}
\titlespacing*{\subsection}{0pt}{5pt plus 1pt minus 1pt}{2pt}

% -----------------------------------------------------------------------------
% Figuras, tabelas, matemática e listas
% -----------------------------------------------------------------------------
\usepackage{graphicx}
\usepackage{booktabs}
\usepackage{tabularx}
\usepackage{array}
\usepackage{amsmath}
\usepackage{siunitx}
\usepackage{subcaption}
\usepackage{balance}
\usepackage{enumitem}
\setlist[itemize]{leftmargin=1.25em,itemsep=0pt,topsep=2pt,parsep=0pt}
\setlist[enumerate]{leftmargin=1.5em,itemsep=0pt,topsep=2pt,parsep=0pt}

% -----------------------------------------------------------------------------
% Legendas compactas
% -----------------------------------------------------------------------------
\usepackage[font=small,labelfont=bf]{caption}
\captionsetup{skip=3pt}
\setlength{\textfloatsep}{7pt plus 2pt minus 2pt}
\setlength{\floatsep}{6pt plus 2pt minus 2pt}
\setlength{\intextsep}{6pt plus 2pt minus 2pt}

% -----------------------------------------------------------------------------
% Referências e links
% -----------------------------------------------------------------------------
\usepackage[backend=biber,style=authoryear,maxcitenames=2,maxbibnames=20]{biblatex}
\addbibresource{referencias.bib}
% Não usar a data de acesso como ano de publicação de páginas sem data.
\DeclareLabeldate{\field{date}\field{year}\literal{nodate}}
\usepackage[hidelinks]{hyperref}
\usepackage{url}

% -----------------------------------------------------------------------------
% Código-fonte
% -----------------------------------------------------------------------------
\usepackage{xcolor}
\usepackage{listings}
\definecolor{codegray}{gray}{0.45}
\definecolor{codegreen}{rgb}{0.0,0.40,0.0}
\definecolor{codeblue}{rgb}{0.10,0.20,0.65}
\lstset{
  basicstyle=\ttfamily\scriptsize,
  keywordstyle=\color{codeblue}\bfseries,
  commentstyle=\color{codegreen},
  numberstyle=\tiny\color{codegray},
  numbers=left,
  numbersep=4pt,
  frame=single,
  aboveskip=4pt,
  belowskip=4pt,
  breaklines=true,
  showstringspaces=false,
  columns=fullflexible
}

% -----------------------------------------------------------------------------
% Orientações didáticas
% Troque para \mostrarorientacoesfalse antes da entrega final.
% -----------------------------------------------------------------------------
\newif\ifmostrarorientacoes
\mostrarorientacoesfalse

\newcommand{\orientacao}[1]{%
  \ifmostrarorientacoes
    \par\smallskip
    \noindent\fcolorbox{black!35}{black!4}{%
      \parbox{0.94\columnwidth}{\fontsize{7.4}{8.2}\selectfont\textbf{Orientação aos autores:} #1}%
    }%
    \par\smallskip
  \fi
}

% =============================================================================
% DOCUMENTO
% =============================================================================
\begin{document}
\twocolumn[
\begin{center}
  {\sffamily\bfseries\fontsize{14}{16}\selectfont
   Comparação de padrões de amostragem Sobol, LP e MMD em renderização\par}
  \vspace{5pt}
  {\fontsize{9.5}{11}\selectfont
   Allan Pinto Martins \quad Felipe Pereira Madureira\par}
  \vspace{2pt}
  {\fontsize{8.5}{10}\selectfont
   Instituto Federal do Norte de Minas Gerais -- IFNMG, Campus Montes Claros\par
   \texttt{apm14@aluno.ifnmg.edu.br, fpm6@aluno.ifnmg.edu.br}\par}
\end{center}
\vspace{3pt}
\begin{minipage}{0.94\textwidth}
\fontsize{8.5}{9.5}\selectfont
\noindent\textbf{Resumo.}
Este projeto investiga o efeito dos padrões Sobol 2D, LP e MMD no erro de renderização com o mesmo número de amostras por pixel. O código de referência foi reproduzido e Sobol foi integrado à câmera do smallpt. A comparação de qualidade entre os padrões será realizada nas próximas etapas.
\par\vspace{3pt}
\noindent\textbf{Palavras-chave:} Computação Gráfica; amostragem; Sobol; renderização.
\par\vspace{5pt}
\noindent\textbf{Abstract.}
This project investigates the effect of Sobol 2D, LP, and MMD sampling patterns on rendering error under equal samples-per-pixel budgets. The reference code has been reproduced, and Sobol has been integrated into camera sampling in smallpt. Image-quality comparisons will be conducted in the next stages.
\par\vspace{3pt}
\noindent\textbf{Keywords:} Computer Graphics; sampling; Sobol; rendering.
\end{minipage}
\vspace{8pt}
]

% Esqueleto para o checkpoint 2. Desenvolver as seções na versão final.
\section{Introdução}
A renderização estima a cor dos pixels usando um número limitado de amostras. A distribuição dessas amostras pode afetar o erro da imagem. A pesquisa parte dos geradores apresentados por \textcite{ahmed2024} e busca responder: com o mesmo número de amostras por pixel, LP e MMD reduzem o erro ou alteram a velocidade de convergência em relação à Sobol 2D convencional?

\section{Trabalho de Referência e Trabalhos Relacionados}
\textcite{ahmed2024} apresenta uma implementação eficiente de Sobol 2D baseada na fatoração da matriz de Pascal, além de procedimentos relacionados aos padrões LP e MMD. Os geradores serão reaproveitados do suplemento oficial. Esta seção será ampliada com a explicação da técnica e a literatura necessária para contextualizar a comparação.

\section{Proposta}
Integrar Sobol, LP e MMD ao mesmo renderizador e comparar o erro das imagens sob condições equivalentes. Sobol será a baseline, ou método de referência. A contribuição da equipe é a integração e a avaliação experimental, mantendo os geradores identificados como código do autor. A alteração será restrita ao par de coordenadas usado na amostragem da câmera.

\section{Implementação}
O suplemento foi compilado em C++ no Windows, pelo MSYS2 MSYS. Sobol foi integrado à câmera do smallpt, de \textcite{smallpt}, preservando a cena e o transporte de luz. A adaptação mantém os quatro subpixels e o filtro tent, aplica deslocamento digital por XOR e separa os estados da câmera e do transporte. LP e MMD ainda não foram integrados.

O código e os resultados estão no repositório\footnote{\url{https://github.com/Felipe-Dev-MP/sobol-rendering-analysis}}, com a primeira integração registrada no commit \texttt{71f0945}. Os detalhes de ambiente e reprodução serão completados na versão final.

\section{Metodologia Experimental}
\subsection{Baseline, variantes e cena}
Serão comparados Sobol, LP e MMD na cena do smallpt, mantendo resolução, câmera, materiais, iluminação e orçamento de amostras iguais. As quantidades de amostras e as sementes das repetições serão definidas após um piloto.

\subsection{Métricas e protocolo}
O erro será medido por RMSE em relação à mesma imagem de referência de alta amostragem. Antes dos experimentos, será preparada uma saída em radiância linear, sem correção gama e sem recorte de cor, e verificada a estabilidade da referência. Serão realizadas repetições com sementes registradas. O protocolo completo ainda será desenvolvido.

\section{Resultados Preliminares}
No teste do suplemento, os modos \texttt{i2x2y}, \texttt{lp}, \texttt{mmd} e \texttt{pbrt} produziram, cada um, 64 pontos distintos em $[0,1)^2$. Os arquivos de \texttt{i2x2y} e \texttt{pbrt} foram idênticos, inclusive na ordem.

A integração Sobol gerou uma imagem de $320\times240$ pixels, com 16 amostras por pixel e semente 1, registrada em \texttt{sobol\_16.ppm}. Isso confirma o funcionamento inicial. Ainda não há comparação de qualidade, referência ou RMSE. As figuras e os resultados quantitativos serão incluídos após os experimentos.

\section{Discussão}
A compilação inicial no UCRT64 falhou por indisponibilidade de \texttt{drand48}; a execução foi viabilizada no MSYS2 MSYS. A taxa negativa de geração impressa no teste foi descartada como medida de desempenho. A análise futura discutirá os resultados e a limitação de alterar apenas a amostragem da câmera, pois outras escolhas do transporte também produzem ruído.

\section{Conclusão}
\textit{A desenvolver após os experimentos.} A versão final responderá à pergunta da pesquisa com base nas medições. Os próximos passos são integrar LP e MMD, preparar a saída linear e executar o piloto com referência, RMSE e repetições.

\balance
\printbibliography[title={Referências}]
\end{document}
